\section{Preliminaries}\label{section: preliminaries}

For $n \in \nat$ we write $\range{n} \defeq \set{1, \dots, n}$.
For a finite linear order $(V, <)$ and $x \le y$, let $\range{x, y} \defeq \set{z \mid x \le z \le y}$.
We write $\succrel$ for the successor relation of $(V, <)$.

\subparagraph*{Terms and $\RSUB$.}
We fix the variables $\vr, \vc, \vb$.
(The lower bound also holds for terms over any finite set of variables, see \Cref{remark: more variables}.)
\emph{Terms} are generated by
\[
  E, F \;::=\; x \mid \emp \mid \id \mid E \union F \mid \com{E} \mid E \compo F \mid E^{*}
  \qquad (x \in \set{\vr, \vc, \vb}).
\]
A \emph{valuation} $\A$ consists of a finite nonempty linear order $(V, <)$ and relations $\vr^{\A}, \vc^{\A}, \vb^{\A}$ contained in the set $\topr_V \defeq \set{(x, y) \in V^2 \mid x \le y}$ of \emph{forward pairs}.
It is an \emph{$\RSUB$ valuation} if each $x^{\A}$ contains either all pairs $(z, z)$ or none of them~\cite{nakamuraFiniteRelationalSemantics2025}.
The value $\sem{E}{\A} \subseteq \topr_V$ is defined as usual.
The constant $\id$ denotes the identity, $\union$ union, $\compo$ relational composition, and ${}^{*}$ reflexive transitive closure.
Complement is taken relative to $\topr_V$: $\sem{\com{E}}{\A} \defeq \topr_V \setminus \sem{E}{\A}$.
Every value consists of forward pairs.
We write $E(x, y)$ for $(x, y) \in \sem{E}{\A}$ when $\A$ is clear.
We use the abbreviations
\[
  E \inter F \defeq \com{\com{E} \union \com{F}},\quad
  E \setminus F \defeq E \inter \com{F},\quad
  \topr \defeq \com{\emp},\quad
  \Dt \defeq \com{\id},\quad
  \St \defeq \Dt \setminus (\Dt \compo \Dt),
\]
and $E^{2} \defeq E \compo E$.
Then $\Dt$ denotes the strict order $<$ and $\St$ denotes $\succrel$.
All of these abbreviations are star-free.

For $x \le y$, let $\A\range{x, y}$ be the valuation induced on $\range{x, y}$.
Every intermediate point of a composition, and every point on a closure path, between $x$ and $y$ lies in $\range{x, y}$.
An induction on terms gives:

\begin{lemma}[Interval locality]\label{lemma: locality}
  For every term $E$, valuation $\A$ and $x \le y$: $(x, y) \in \sem{E}{\A}$ iff $(x, y) \in \sem{E}{\A\range{x, y}}$.
\end{lemma}

\subparagraph*{Star height.}
The \emph{star height} $\sh(E)$ is the nesting depth of ${}^{*}$ in $E$.
Complement and the other operators do not count.
So $\sh(x) = \sh(\emp) = \sh(\id) = 0$, $\sh(\com{E}) = \sh(E)$, $\sh(E \union F) = \sh(E \compo F) = \max(\sh(E), \sh(F))$, and $\sh(E^*) = \sh(E) + 1$.
For a class $\mathcal{C}$ of valuations, $E$ and $F$ are \emph{equivalent on $\mathcal{C}$} if $\sem{E}{\A} = \sem{F}{\A}$ for all $\A \in \mathcal{C}$.
The \emph{generalized star height of $E$ relative to $\mathcal{C}$}, written $\gsh{\mathcal{C}}(E)$, is the least $\sh(F)$ over all $F$ equivalent to $E$ on $\mathcal{C}$.

\subparagraph*{Pathwidth.}
A \emph{path decomposition} of a graph $G$ is a sequence of \emph{bags} $X_1, \dots, X_\ell \subseteq V(G)$ with the following properties.
Every edge is contained in some bag, and for every vertex, the bags containing it form a nonempty interval.
Its width is $\max_i \card{X_i} - 1$, and the \emph{pathwidth} $\pw(G)$ of $G$ is the least width of a path decomposition of $G$~\cite{robertsonGraphMinorsExcluding1983}.
Forests have treewidth at most one, and every graph with a cycle has treewidth, and hence pathwidth, at least two.

The \emph{graph} $\Gaif(\A)$ of a valuation $\A$ has vertex set $V$.
Two distinct vertices $x, y$ are adjacent iff $(x, y)$ or $(y, x)$ belongs to $\succrel \union \vr^{\A} \union \vc^{\A} \union \vb^{\A}$.
So $\Gaif(\A)$ is the Gaifman graph of $(V, \succrel, \vr^{\A}, \vc^{\A}, \vb^{\A})$ with loops removed.
Only input relations and the successor contribute edges, and derived relations such as $\sem{\vr \compo \vr}{\A}$ do not.
We write $\pw(\A) \defeq \pw(\Gaif(\A))$ and
\[
  \Cls{k} \defeq \set{\A \mid \A \text{ is an $\RSUB$ valuation with } \pw(\A) \le k}.
\]

\subparagraph*{Star-free terms and first-order types.}
A term is \emph{star-free} if it contains no ${}^{*}$.
The \emph{composition depth} $\cd(E)$ of a star-free term is the nesting depth of $\compo$.
Thus $\cd$ is $0$ on atoms, is preserved by $\union$ and complement, and $\cd(E \compo F) = \max(\cd(E), \cd(F)) + 1$.
Let $\SF{q}$ be the set of star-free terms $E$ with $\cd(E) \le q$.
Since $\Dt$ has depth $0$, $E \in \SF{q}$ implies $E \inter \Dt \in \SF{q}$.

We regard a valuation as a first-order structure over the signature $\set{<, \vr, \vc, \vb}$.
For a structure $\mathfrak{B}$ and elements $\bar{a}$, $\tp{q}(\mathfrak{B}, \bar{a})$ denotes the first-order type of quantifier rank $q$ of $\bar a$ in $\mathfrak{B}$.
We use the Ehrenfeucht--Fra\"iss\'e characterization of these types~\cite{ehrenfeuchtApplicationGamesCompleteness1961,fraisseClassificationsSystemsRelations1954,ebbinghausFiniteModelTheory1995}.
The standard translation sends $E \compo F$ to $\exists z\,(x \le z \le y \wedge \varphi_E(x, z) \wedge \varphi_F(z, y))$.
It turns each $E \in \SF{q}$ into a first-order formula $\varphi_E(x, y)$ of quantifier rank at most $q$ with $\sem{E}{\A} = \set{(x, y) \in \topr_V \mid \A \models \varphi_E(x, y)}$.
With \Cref{lemma: locality} this gives:

\begin{lemma}\label{lemma: SF types}
  For $E \in \SF{q}$ and $x \le y$, whether $(x, y) \in \sem{E}{\A}$ holds depends only on the type $\tp{q}(\A\range{x, y}, x, y)$.
  Moreover, $\SF{q}$ contains only finitely many terms up to equivalence on all valuations.
\end{lemma}

We fix a finite set $\Rep{q} \subseteq \SF{q}$ of representatives of $\SF{q}$ up to equivalence.

\subparagraph*{Words and their types.}
Let $\Gamma$ be a finite alphabet.
A nonempty word $\alpha = \alpha_1 \cdots \alpha_k \in \Gamma^{+}$ is the structure with universe $\range{k}$, the order $<$, the successor, a unary predicate for each letter, and the constants $1$ and $k$.
Let $\tau_q(\alpha)$ be its first-order type of quantifier rank $q$.
For the empty word we add a fresh value $\tau_q(\varepsilon)$.
The following facts are standard.
We include proofs, since we need the exact thresholds.

\begin{lemma}\label{lemma: word types}
  Let $q \in \nat$.
  \begin{enumerate}
    \item\label{item: composition} If $\tau_q(\alpha) = \tau_q(\alpha')$ and $\tau_q(\beta) = \tau_q(\beta')$, then $\tau_q(\alpha\beta) = \tau_q(\alpha'\beta')$.
          Hence the values of $\tau_q$ form a finite monoid $\Types{q}(\Gamma)$, and $\tau_q \colon \Gamma^{*} \to \Types{q}(\Gamma)$ is a surjective monoid morphism.
    \item\label{item: powers} For every $\omega \in \Gamma^{+}$ and all $k, k' > 2^{q+1}$, we have $\tau_q(\omega^{k}) = \tau_q(\omega^{k'})$.
          In particular, $\Types{q}(\Gamma)$ is aperiodic.
  \end{enumerate}
\end{lemma}
\begin{proof}
  (\ref{item: composition})
  Duplicator plays the two games on $\alpha, \alpha'$ and on $\beta, \beta'$ in parallel.
  She answers each move in the part where it was made.
  The constants are matched, so the last position of $\alpha$ corresponds to the last position of $\alpha'$, and similarly for the first positions of $\beta, \beta'$.
  The only successor pair across the two parts joins the last position of $\alpha$ and the first position of $\beta$.
  So the union of the two partial isomorphisms is a partial isomorphism.
  The empty word is the identity.

  (\ref{item: powers})
  Write the positions of $\omega^{k}$ as pairs $(j, p)$, where $j \in \range{0, k-1}$ is the index of the copy and $p \in \range{\card{\omega}}$ is the position inside it.
  Duplicator always answers with the same $p$.
  She plays on the copy indices the usual strategy for linear orders~\cite{ebbinghausFiniteModelTheory1995}.
  This strategy keeps the following invariant for all pebbled elements, including the two constants.
  When $r$ rounds remain, any two copy indices have the same difference on both sides, or both differences are at least $2^{r+1}$ in absolute value, with the same sign.
  The invariant holds initially because $k-1, k'-1 \ge 2^{q+1}$.
  At the end, copy indices have equal differences or differ by at least $2$.
  Hence order, letters and the constants are preserved.
  The successor $(j, p) \to (j', p')$ requires $j' - j \in \set{0, 1}$, so it is preserved too.
  Taking $\omega$ to be a representative of an element $m \in \Types{q}(\Gamma)$ gives $m^{k} = m^{k+1}$ for $k > 2^{q+1}$.
\end{proof}

\subparagraph*{Finite monoids.}
Every finite monoid $\Tm$ has a least nonempty two-sided ideal $\Kid(\Tm)$, its \emph{minimal ideal}, and $\Kid(\Tm)$ contains an idempotent.
A monoid is \emph{aperiodic} if all its subgroups are trivial.

\begin{lemma}\label{lemma: minimal ideal}
  Let $\Tm$ be a finite monoid and $e \in \Kid(\Tm)$ an idempotent.
  \begin{enumerate}
    \item\label{item: group} $e \Tm e$ is a group with identity $e$.
    \item\label{item: collapse} If $\pi \colon \Tm \to \Pm$ is a surjective morphism onto an aperiodic monoid, then $\pi(e)\, n\, \pi(e) = \pi(e)$ for all $n \in \Pm$.
  \end{enumerate}
\end{lemma}
\begin{proof}
  (\ref{item: group})
  Let $x \in e\Tm e$.
  Since $\Tm x \Tm$ is an ideal contained in $\Kid(\Tm)$, it equals $\Kid(\Tm)$, so $e = a x b$ for some $a, b$.
  Then $e = (eae)\, x\, (ebe)$.
  In the finite monoid $e\Tm e$, a one-sided inverse is a two-sided inverse.
  So $(eae)x$ and $x(ebe)$ are invertible, and hence $x$ is invertible.

  (\ref{item: collapse})
  Since $\pi$ is surjective, $\pi(\Kid(\Tm))$ is an ideal of $\Pm$.
  Every ideal $J$ of $\Pm$ satisfies $\pi^{-1}(J) \supseteq \Kid(\Tm)$.
  So $\pi(\Kid(\Tm)) = \Kid(\Pm)$, and $\pi(e)$ is an idempotent in $\Kid(\Pm)$.
  By (\ref{item: group}), $\pi(e)\Pm\pi(e)$ is a subgroup of $\Pm$, so it is trivial.
\end{proof}
